% ============================================================
%  C: Como Programar -- 6ª Edição | Deitel & Deitel
%  Capítulo 14 -- Outros Tópicos sobre C
%  EXERCÍCIOS (14.2 -- 14.7)
%  Abordagem incremental: Caps. 1 a 14
% ============================================================
\documentclass[12pt,a4paper]{article}
\usepackage[utf8]{inputenc}
\usepackage[T1]{fontenc}
\usepackage[brazil]{babel}
\usepackage{geometry}
\geometry{top=2.5cm,bottom=2.5cm,left=3cm,right=3cm}
\usepackage{parskip}\setlength{\parskip}{6pt}\setlength{\parindent}{0pt}
\usepackage{xcolor,tcolorbox,enumitem,listings,fancyhdr,titlesec,hyperref,booktabs,array,amsmath}
\tcbuselibrary{skins,breakable}

\definecolor{deitelblue}{RGB}{0,70,127}
\definecolor{deitelgray}{RGB}{240,242,245}
\definecolor{deitelgreen}{RGB}{0,110,60}
\definecolor{deitellightblue}{RGB}{220,235,250}
\definecolor{deitelred}{RGB}{160,20,20}
\definecolor{deitellorange}{RGB}{200,90,0}

\newtcolorbox{boaprática}[1][]{colback=deitellightblue,colframe=deitelblue,
  fonttitle=\bfseries\small,title={Boa prática de programação},breakable,#1}
\newtcolorbox{erroco}[1][]{colback=red!5,colframe=deitelred,
  fonttitle=\bfseries\small,title={Erro comum de programação},breakable,#1}
\newtcolorbox{respostabox}[1][]{colback=deitelgray,colframe=deitelblue!60,
  fonttitle=\bfseries\small\color{deitelblue},title={Resposta detalhada},breakable,#1}
\newtcolorbox{enunciadoBox}[1][]{colback=white,colframe=deitelblue,
  fonttitle=\bfseries,title={#1},breakable}
\newtcolorbox{saida}[1][]{colback=black!88,colframe=deitelblue,
  fontupper=\ttfamily\small\color{green!70},
  title={\small\color{white}Saída do programa},breakable,#1}
\newtcolorbox{atencao}[1][]{colback=yellow!10,colframe=deitellorange,
  fonttitle=\bfseries\small,title={Atenção},breakable,#1}

\setlist[enumerate,1]{label=\textbf{\alph*)},leftmargin=2em}
\setlist[itemize]{leftmargin=2em}

\lstset{
  basicstyle=\ttfamily\small,backgroundcolor=\color{deitelgray},
  frame=single,framerule=0.4pt,rulecolor=\color{deitelblue!40},
  breaklines=true,showstringspaces=false,
  keywordstyle=\color{deitelblue}\bfseries,
  commentstyle=\color{deitelgreen}\itshape,
  stringstyle=\color{deitelred},
  language=C,numbers=left,numberstyle=\tiny\color{gray},
  xleftmargin=18pt,xrightmargin=5pt,stepnumber=1,
}

\pagestyle{fancy}\fancyhf{}
\fancyhead[L]{\small\color{deitelblue}\textbf{C: Como Programar -- 6ª ed. | Deitel \& Deitel}}
\fancyhead[R]{\small\color{deitelblue}Cap. 14 -- Exercícios}
\fancyfoot[C]{\small\thepage}
\renewcommand{\headrulewidth}{0.4pt}

\titleformat{\section}[block]{\large\bfseries\color{deitelblue}}{\thesection.}{0.5em}{}[\titlerule]
\titleformat{\subsection}[block]{\normalsize\bfseries\color{deitelblue!80}}{}{}{}

\hypersetup{colorlinks=true,linkcolor=deitelblue}

\begin{document}

\begin{titlepage}
  \centering\vspace*{2cm}
  {\color{deitelblue}\rule{\linewidth}{2pt}}\\[0.4cm]
  {\huge\bfseries\color{deitelblue}C: Como Programar}\\[0.2cm]
  {\large\color{deitelblue}6ª Edição $\cdot$ Paul Deitel \& Harvey Deitel}\\[0.2cm]
  {\color{deitelblue}\rule{\linewidth}{2pt}}\\[1cm]
  {\LARGE\bfseries Capítulo 14}\\[0.3cm]
  {\Large\color{deitelblue!80}Outros Tópicos sobre C}\\[0.8cm]
  \begin{tcolorbox}[colback=deitellightblue,colframe=deitelblue,width=0.85\linewidth]
    \centering{\large\bfseries Exercícios (14.2--14.7)}\\[0.2cm]
    {\normalsize Resolvidos passo a passo, com código completo\\[0.2cm]
    \textbf{O capítulo final do núcleo de C}}
  \end{tcolorbox}
\end{titlepage}

\tableofcontents\newpage

% ============================================================
\section{Exercício 14.2 -- Produto com Argumentos Variáveis}
% ============================================================

\begin{enunciadoBox}[Enunciado 14.2]
Escreva um programa que calcule o produto de uma série de inteiros passados à
função \texttt{product} por uma lista de argumentos de tamanhos variáveis.
Teste com várias chamadas, cada uma com um número diferente de argumentos.
\end{enunciadoBox}

\begin{respostabox}
\begin{lstlisting}
/* Ex14.2: produto_variadic.c
   Demonstra uso de stdarg.h para funcao variadica */
#include <stdio.h>
#include <stdarg.h>

/* O primeiro parametro fixo eh a quantidade de inteiros que segue */
int product( int contagem, ... )
{
    va_list args;
    int i, prod = 1;

    va_start( args, contagem );

    for ( i = 0; i < contagem; ++i ) {
        prod *= va_arg( args, int );
    }

    va_end( args );  /* SEMPRE chamar antes do return! */
    return prod;
}

int main( void )
{
    /* 2 argumentos */
    printf( "Produto de 2 e 3:        %d\n",
            product( 2, 2, 3 ) );

    /* 3 argumentos */
    printf( "Produto de 4, 5 e 6:     %d\n",
            product( 3, 4, 5, 6 ) );

    /* 5 argumentos */
    printf( "Produto de 1, 2, 3, 4, 5: %d\n",
            product( 5, 1, 2, 3, 4, 5 ) );

    /* 1 argumento */
    printf( "Produto de 7:            %d\n",
            product( 1, 7 ) );

    /* nenhum argumento variavel (so o contador) */
    printf( "Produto vazio:           %d\n",
            product( 0 ) );

    return 0;
}
\end{lstlisting}

\textbf{Saída:}
\begin{saida}
Produto de 2 e 3:        6\\
Produto de 4, 5 e 6:     120\\
Produto de 1, 2, 3, 4, 5: 120\\
Produto de 7:            7\\
Produto vazio:           1
\end{saida}
\end{respostabox}

\subsection*{Análise didática}

\begin{boaprática}
\textbf{Anatomia de uma função variádica:}

\begin{enumerate}[noitemsep]
  \item \textbf{Pelo menos um parâmetro nomeado} -- aqui \texttt{contagem}. As
  reticências (\texttt{...}) sempre vêm \textit{depois} de pelo menos um parâmetro
  fixo, que serve como ``âncora''.

  \item \textbf{Declarar \texttt{va\_list}} -- tipo opaco que representa o estado
  do percurso pelos argumentos.

  \item \textbf{\texttt{va\_start}} -- inicializa, recebendo como segundo
  argumento o \textit{último parâmetro nomeado} (não as reticências).

  \item \textbf{Iterar com \texttt{va\_arg}} -- para cada argumento, especificamos
  o tipo esperado. Aqui sempre \texttt{int}.

  \item \textbf{\texttt{va\_end}} -- liberar recursos internos. Não chamar pode
  causar comportamento indefinido.
\end{enumerate}
\end{boaprática}

\bigskip

\begin{erroco}
\textbf{Armadilhas de funções variádicas:}

\begin{itemize}[noitemsep]
  \item \textbf{Sem checagem de tipos:} \texttt{product(2, 3.14, 5)} compila mas
  causa estrago -- \texttt{va\_arg(args, int)} interpreta os bytes do \texttt{double}
  como \texttt{int}.

  \item \textbf{Promoções automáticas:} \texttt{char} e \texttt{short} são promovidos
  a \texttt{int}; \texttt{float} a \texttt{double}. Nunca use \texttt{va\_arg(args, char)}
  ou \texttt{va\_arg(args, float)} -- comportamento indefinido.

  \item \textbf{Precisa de sentinela ou contador:} a função não sabe quantos
  argumentos foram passados. Por isso passamos \texttt{contagem} explicitamente.
  \texttt{printf} sabe pelo número de \texttt{\%} no formato; outras funções
  usam um valor sentinela como \texttt{NULL} ou \texttt{-1}.
\end{itemize}
\end{erroco}

\newpage

% ============================================================
\section{Exercício 14.3 -- Argumentos da Linha de Comando}
% ============================================================

\begin{enunciadoBox}[Enunciado 14.3]
Escreva um programa que imprima os argumentos da linha de comando.
\end{enunciadoBox}

\begin{respostabox}
\begin{lstlisting}
/* Ex14.3: argumentos.c */
#include <stdio.h>

int main( int argc, char *argv[] )
{
    int i;

    printf( "Numero de argumentos (argc): %d\n\n", argc );

    printf( "Argumentos:\n" );
    for ( i = 0; i < argc; ++i ) {
        printf( "  argv[%d] = \"%s\"\n", i, argv[ i ] );
    }

    return 0;
}
\end{lstlisting}

\textbf{Compilação e execução:}
\begin{lstlisting}
$ gcc -o argumentos argumentos.c
$ ./argumentos ola mundo arquivo.txt 42
\end{lstlisting}

\textbf{Saída:}
\begin{saida}
Numero de argumentos (argc): 5\\
\\
Argumentos:\\
\ \ argv[0] = "./argumentos"\\
\ \ argv[1] = "ola"\\
\ \ argv[2] = "mundo"\\
\ \ argv[3] = "arquivo.txt"\\
\ \ argv[4] = "42"
\end{saida}
\end{respostabox}

\subsection*{Análise didática}

\begin{atencao}
\textbf{Pontos cruciais sobre \texttt{argc}/\texttt{argv}:}

\begin{itemize}[noitemsep]
  \item \textbf{\texttt{argv[0]} é o nome do programa} (geralmente, com o caminho).
  Em alguns sistemas pode ser o caminho completo, em outros apenas o nome.

  \item \textbf{\texttt{argc} inclui o nome do programa.} Por isso \texttt{argc \(\geq\) 1}
  sempre (em sistemas conformes).

  \item \textbf{\texttt{argv[argc]} é garantidamente \texttt{NULL}.} Você pode
  iterar até encontrar \texttt{NULL} em vez de usar \texttt{argc}.

  \item \textbf{Tudo é string.} Para usar \texttt{argv[1]} como número, converta
  com \texttt{atoi}, \texttt{atof} ou \texttt{strtol}.

  \item \textbf{O shell faz o ``split''.} \texttt{./prog "ola mundo"} resulta em
  \texttt{argc = 2}, \texttt{argv[1] = "ola mundo"}. Sem aspas, seriam dois
  argumentos separados.
\end{itemize}

\textbf{Forma alternativa equivalente:}
\begin{lstlisting}
int main( int argc, char **argv )  /* ponteiro para ponteiro */
\end{lstlisting}
As duas formas (\texttt{char *argv[]} e \texttt{char **argv}) são intercambiáveis
para parâmetros de função.
\end{atencao}

\newpage

% ============================================================
\section{Exercício 14.4 -- Ordenar com -c ou -d}
% ============================================================

\begin{enunciadoBox}[Enunciado 14.4]
Escreva um programa que classifique um array de inteiros em ordem crescente
ou decrescente. O programa deve usar argumentos da linha de comando para
passar \texttt{-c} (crescente) ou \texttt{-d} (decrescente).
\end{enunciadoBox}

\begin{respostabox}
\begin{lstlisting}
/* Ex14.4: ordena_cli.c
   Uso:
     ./ordena -c 5 3 8 1 4    (crescente)
     ./ordena -d 5 3 8 1 4    (decrescente) */
#include <stdio.h>
#include <stdlib.h>
#include <string.h>

void bubbleSort( int *arr, int n, int crescente )
{
    int i, j, temp;
    for ( i = 0; i < n - 1; ++i ) {
        for ( j = 0; j < n - 1 - i; ++j ) {
            int trocar = crescente ? ( arr[ j ] > arr[ j + 1 ] )
                                    : ( arr[ j ] < arr[ j + 1 ] );
            if ( trocar ) {
                temp         = arr[ j ];
                arr[ j ]     = arr[ j + 1 ];
                arr[ j + 1 ] = temp;
            }
        }
    }
}

int main( int argc, char *argv[] )
{
    int i, n, crescente;

    if ( argc < 3 ) {
        printf( "Uso: %s -c|-d num1 num2 ...\n", argv[ 0 ] );
        return 1;
    }

    /* Determina ordem */
    if ( strcmp( argv[ 1 ], "-c" ) == 0 ) {
        crescente = 1;
    } else if ( strcmp( argv[ 1 ], "-d" ) == 0 ) {
        crescente = 0;
    } else {
        printf( "Opcao invalida: %s. Use -c ou -d.\n", argv[ 1 ] );
        return 1;
    }

    /* Quantidade de numeros */
    n = argc - 2;

    /* Aloca array dinamicamente -- conceito do Cap. 12 */
    int *arr = malloc( n * sizeof( int ) );
    if ( arr == NULL ) {
        printf( "Erro: falta de memoria.\n" );
        return 1;
    }

    /* Converte strings em inteiros */
    for ( i = 0; i < n; ++i ) {
        arr[ i ] = atoi( argv[ i + 2 ] );
    }

    /* Ordena */
    bubbleSort( arr, n, crescente );

    /* Exibe resultado */
    printf( "Resultado (%s):\n",
            crescente ? "crescente" : "decrescente" );
    for ( i = 0; i < n; ++i ) {
        printf( "%d ", arr[ i ] );
    }
    printf( "\n" );

    free( arr );
    return 0;
}
\end{lstlisting}

\textbf{Exemplos de uso:}
\begin{saida}
\$ ./ordena -c 5 3 8 1 4 9 2\\
Resultado (crescente):\\
1 2 3 4 5 8 9\\
\\
\$ ./ordena -d 5 3 8 1 4 9 2\\
Resultado (decrescente):\\
9 8 5 4 3 2 1
\end{saida}
\end{respostabox}

\subsection*{Análise didática}

\begin{boaprática}
\textbf{Padrão Unix de opções:} flags começam com \texttt{-} seguidas de letra
(\texttt{-c}, \texttt{-d}, \texttt{-v}, etc.). Para múltiplas opções, há
convenções:

\begin{itemize}[noitemsep]
  \item \textbf{Forma curta:} \texttt{-c}, \texttt{-d}, \texttt{-v}.
  \item \textbf{Forma longa:} \texttt{--crescente}, \texttt{--decrescente},
  \texttt{--verbose} (GNU).
  \item \textbf{Combináveis:} \texttt{-cvd} equivale a \texttt{-c -v -d}.
  \item \textbf{Com valor:} \texttt{-o saida.txt} ou \texttt{--output=saida.txt}.
\end{itemize}

Para projetos sérios, use a biblioteca \texttt{<unistd.h>} com \texttt{getopt}
(POSIX) ou \texttt{getopt\_long} (GNU). Aqui implementamos manualmente para
mostrar o conceito.
\end{boaprática}

\newpage

% ============================================================
\section{Exercício 14.5 -- Arquivos Temporários}
% ============================================================

\begin{enunciadoBox}[Enunciado 14.5]
Escreva um programa que coloque um espaço entre cada caractere em um arquivo.
Use um arquivo temporário para isso: primeiro grave o conteúdo modificado no
temporário, depois copie de volta ao arquivo original.
\end{enunciadoBox}

\begin{respostabox}
\begin{lstlisting}
/* Ex14.5: espaca_chars.c
   Insere espaco entre cada caractere de um arquivo */
#include <stdio.h>
#include <stdlib.h>

int main( int argc, char *argv[] )
{
    FILE *original, *temp;
    int c;

    if ( argc != 2 ) {
        printf( "Uso: %s arquivo.txt\n", argv[ 0 ] );
        return 1;
    }

    /* Abre o original para leitura */
    if ( ( original = fopen( argv[ 1 ], "r" ) ) == NULL ) {
        printf( "Erro: nao foi possivel abrir %s\n", argv[ 1 ] );
        return 1;
    }

    /* Cria arquivo temporario */
    if ( ( temp = tmpfile() ) == NULL ) {
        printf( "Erro: nao foi possivel criar arquivo temporario.\n" );
        fclose( original );
        return 1;
    }

    /* Etapa 1: copia do original para o temporario inserindo espacos */
    while ( ( c = fgetc( original ) ) != EOF ) {
        if ( c != '\n' && c != ' ' ) {
            fputc( c, temp );
            fputc( ' ', temp );  /* espaco apos cada caractere */
        } else {
            fputc( c, temp );    /* preserva quebras e espacos */
        }
    }
    fclose( original );

    /* Etapa 2: rebobina temporario e reabre original para escrita */
    rewind( temp );
    if ( ( original = fopen( argv[ 1 ], "w" ) ) == NULL ) {
        printf( "Erro: nao foi possivel reescrever %s\n", argv[ 1 ] );
        fclose( temp );
        return 1;
    }

    /* Etapa 3: copia do temporario para o original */
    while ( ( c = fgetc( temp ) ) != EOF ) {
        fputc( c, original );
    }

    fclose( original );
    fclose( temp );  /* tmpfile() eh deletado automaticamente */

    printf( "Arquivo %s processado com sucesso.\n", argv[ 1 ] );
    return 0;
}
\end{lstlisting}

\textbf{Exemplo:}
\begin{saida}
\$ cat exemplo.txt\\
Ola mundo\\
\\
\$ ./espaca\_chars exemplo.txt\\
Arquivo exemplo.txt processado com sucesso.\\
\\
\$ cat exemplo.txt\\
O l a  m u n d o
\end{saida}

\textbf{Vantagens de \texttt{tmpfile()}:}
\begin{itemize}[noitemsep]
  \item Nome único garantido pelo sistema.
  \item \textbf{Arquivo deletado automaticamente} ao fechar ou ao terminar o programa.
  \item Sem risco de colisão de nomes com outros processos.
  \item Aberto em modo binário \texttt{"wb+"} (leitura e escrita).
\end{itemize}
\end{respostabox}

\newpage

% ============================================================
\section{Exercício 14.6 -- Tratamento de Sinais}
% ============================================================

\begin{enunciadoBox}[Enunciado 14.6]
Escreva um programa com tratadores de sinal para \texttt{SIGABRT} e
\texttt{SIGINT}. Teste chamando \texttt{abort} (gera SIGABRT) e digitando
\texttt{Ctrl+C} (gera SIGINT).
\end{enunciadoBox}

\begin{respostabox}
\begin{lstlisting}
/* Ex14.6: sinais.c */
#include <stdio.h>
#include <stdlib.h>
#include <signal.h>

void tratadorSinal( int sinalRecebido )
{
    /* IMPORTANTE: re-registrar o handler porque alguns sistemas
       resetam para SIG_DFL apos receber o sinal               */
    signal( sinalRecebido, tratadorSinal );

    printf( "\n*** Sinal capturado: " );

    switch ( sinalRecebido ) {
        case SIGABRT:
            printf( "SIGABRT (abort) ***\n" );
            printf( "    Programa solicitou aborto.\n" );
            break;
        case SIGINT:
            printf( "SIGINT (Ctrl+C) ***\n" );
            printf( "    Interrupcao do usuario detectada.\n" );
            break;
        default:
            printf( "Sinal %d ***\n", sinalRecebido );
    }

    printf( "    Encerrando programa de forma controlada.\n" );
    exit( EXIT_SUCCESS );
}

int main( void )
{
    int opcao;

    /* Registra tratadores */
    signal( SIGABRT, tratadorSinal );
    signal( SIGINT,  tratadorSinal );

    printf( "Programa de teste de sinais\n" );
    printf( "  1 - gerar SIGABRT (via abort)\n" );
    printf( "  2 - gerar SIGABRT (via raise)\n" );
    printf( "  3 - gerar SIGINT (via raise)\n" );
    printf( "  4 - loop infinito (use Ctrl+C para testar SIGINT)\n" );
    printf( "Escolha: " );
    scanf( "%d", &opcao );

    switch ( opcao ) {
        case 1: abort();              break;
        case 2: raise( SIGABRT );     break;
        case 3: raise( SIGINT );      break;
        case 4:
            printf( "Pressione Ctrl+C para interromper...\n" );
            for ( ;; ) ;   /* loop infinito */
            break;
        default:
            printf( "Opcao invalida.\n" );
    }

    return 0;
}
\end{lstlisting}

\textbf{Exemplo de execução:}
\begin{saida}
\$ ./sinais\\
Programa de teste de sinais\\
\ \ 1 - gerar SIGABRT (via abort)\\
\ \ 2 - gerar SIGABRT (via raise)\\
\ \ 3 - gerar SIGINT (via raise)\\
\ \ 4 - loop infinito\\
Escolha: 4\\
Pressione Ctrl+C para interromper...\\
\^{}C\\
*** Sinal capturado: SIGINT (Ctrl+C) ***\\
\ \ \ \ Interrupcao do usuario detectada.\\
\ \ \ \ Encerrando programa de forma controlada.
\end{saida}
\end{respostabox}

\subsection*{Sinais comuns em sistemas Unix-like}

\begin{respostabox}
\begin{center}
\renewcommand{\arraystretch}{1.25}
\begin{tabular}{l l l}
\toprule
\textbf{Sinal} & \textbf{Significado} & \textbf{Causa típica} \\ \midrule
\texttt{SIGABRT} & Aborto                & chamada a \texttt{abort()} \\
\texttt{SIGFPE}  & Erro aritmético        & divisão por zero \\
\texttt{SIGILL}  & Instrução ilegal       & corrupção de código \\
\texttt{SIGINT}  & Interrupção            & Ctrl+C no terminal \\
\texttt{SIGSEGV} & Violação de segmentação& ponteiro inválido \\
\texttt{SIGTERM} & Término solicitado     & \texttt{kill <pid>} \\
\bottomrule
\end{tabular}
\end{center}

\begin{atencao}
\textbf{Cuidados com tratadores de sinais:}
\begin{itemize}[noitemsep]
  \item Tratadores devem ser \textbf{curtos e simples}. Idealmente, apenas
  setar uma flag global \texttt{volatile sig\_atomic\_t}.
  \item \textbf{Nem todas as funções são seguras} em tratadores. \texttt{printf}
  e \texttt{malloc} \textit{não são} async-signal-safe; podem deadlock.
  \item \texttt{SIGKILL} e \texttt{SIGSTOP} \textbf{não podem ser interceptados}.
  São o ``mecanismo de força bruta'' do SO.
\end{itemize}
\end{atencao}
\end{respostabox}

\newpage

% ============================================================
\section{Exercício 14.7 -- Alocação Dinâmica com \texttt{realloc}}
% ============================================================

\begin{enunciadoBox}[Enunciado 14.7]
Escreva um programa que aloque um array de inteiros dinamicamente, com tamanho
fornecido pelo teclado. Leia os elementos, imprima-os. Em seguida, redimensione
para a metade do tamanho original e imprima os elementos restantes.
\end{enunciadoBox}

\begin{respostabox}
\begin{lstlisting}
/* Ex14.7: realloc_demo.c
   Demonstra malloc, realloc e free */
#include <stdio.h>
#include <stdlib.h>

void imprimirArray( const int *arr, int n, const char *titulo )
{
    int i;
    printf( "%s (%d elementos):\n  ", titulo, n );
    for ( i = 0; i < n; ++i ) {
        printf( "%d ", arr[ i ] );
    }
    printf( "\n" );
}

int main( void )
{
    int  tamanho, i;
    int *arr;

    printf( "Tamanho do array: " );
    scanf( "%d", &tamanho );

    if ( tamanho <= 0 ) {
        printf( "Tamanho invalido.\n" );
        return 1;
    }

    /* Aloca com calloc (zera tudo) */
    arr = calloc( tamanho, sizeof( int ) );
    if ( arr == NULL ) {
        printf( "Erro: memoria insuficiente.\n" );
        return 1;
    }

    /* Le valores */
    printf( "Digite %d inteiros:\n", tamanho );
    for ( i = 0; i < tamanho; ++i ) {
        scanf( "%d", &arr[ i ] );
    }

    imprimirArray( arr, tamanho, "Array original" );

    /* Realoca para metade do tamanho */
    int novoTamanho = tamanho / 2;
    int *temp = realloc( arr, novoTamanho * sizeof( int ) );

    if ( temp == NULL ) {
        printf( "Erro ao realocar.\n" );
        free( arr );  /* libera bloco original */
        return 1;
    }
    arr = temp;  /* atualiza ponteiro */

    imprimirArray( arr, novoTamanho, "Array apos realloc (metade)" );

    free( arr );
    return 0;
}
\end{lstlisting}

\textbf{Exemplo de execução:}
\begin{saida}
Tamanho do array: 6\\
Digite 6 inteiros:\\
10 20 30 40 50 60\\
Array original (6 elementos):\\
\ \ 10 20 30 40 50 60\\
Array apos realloc (metade) (3 elementos):\\
\ \ 10 20 30
\end{saida}
\end{respostabox}

\subsection*{Anatomia de \texttt{realloc}}

\begin{boaprática}
\textbf{Como funciona \texttt{realloc}:}
\begin{lstlisting}
void *realloc( void *ptr, size_t novoTamanho );
\end{lstlisting}

\textbf{Casos possíveis:}
\begin{enumerate}[noitemsep]
  \item \textbf{Tamanho menor:} mantém os bytes iniciais, libera o resto.
  Pode retornar o mesmo ponteiro.
  \item \textbf{Tamanho maior, espaço contíguo disponível:} expande no lugar.
  \item \textbf{Tamanho maior, sem espaço:} aloca novo bloco em outro lugar,
  \textit{copia} os dados antigos, libera o original. \textbf{O ponteiro muda!}
  \item \textbf{Falha:} retorna \texttt{NULL}, o ponteiro original \textbf{permanece válido}.
\end{enumerate}

\textbf{Por isso o idioma:}
\begin{lstlisting}
int *temp = realloc( arr, novoTam * sizeof( int ) );
if ( temp != NULL ) {
    arr = temp;          /* sucesso */
} else {
    /* falha: arr ainda eh valido, precisa de free dele */
    free( arr );
    return 1;
}
\end{lstlisting}

\begin{erroco}
\textbf{Erro classico:}
\begin{lstlisting}
arr = realloc( arr, novoTam );  /* PERIGOSO */
\end{lstlisting}
Se \texttt{realloc} falhar, \texttt{arr} vira \texttt{NULL} e perdemos
\textit{permanentemente} o ponteiro para a memória original -- \textbf{memory leak}!
\end{erroco}
\end{boaprática}

\newpage

% ============================================================
\section*{Reflexão Final -- Encerrando o Núcleo de C}

\begin{tcolorbox}[colback=deitellightblue,colframe=deitelblue,
  title={\bfseries 14 capítulos, uma jornada completa}]

Chegamos ao fim do núcleo de C. Olhando para trás, completamos uma travessia
notável:

\textbf{Capítulos 1--2 -- Os primeiros passos:} de zero a programas funcionais.

\textbf{Capítulos 3--5 -- Programação estruturada:} controle de fluxo, funções,
recursão.

\textbf{Capítulos 6--8 -- Dados compostos:} arrays, ponteiros (o coração de C!), strings.

\textbf{Capítulo 9 -- E/S formatada:} domínio completo de \texttt{printf}/\texttt{scanf}.

\textbf{Capítulos 10--12 -- Estruturas e persistência:} \texttt{struct},
arquivos, listas, pilhas, filas, árvores.

\textbf{Capítulo 13 -- O pré-processador:} a meta-linguagem do compilador.

\textbf{Capítulo 14 -- Outros tópicos:} variádicas, linha de comando, sinais,
\texttt{realloc}, \texttt{Makefile}.

\bigskip

\textbf{O que você ganhou:}
\begin{itemize}[noitemsep]
  \item Capacidade de escrever programas profissionais em C.
  \item Compreensão de baixo nível: memória, ponteiros, bits, arquivos.
  \item Vocabulário para entender qualquer código C existente.
  \item Base sólida para qualquer linguagem moderna -- C++, Rust, Go, Java
  herdam grande parte da sintaxe e semântica.
  \item Pré-requisito para sistemas operacionais, compiladores, embarcados,
  drivers, jogos, alta performance.
\end{itemize}

\textbf{Próximos passos sugeridos:}
\begin{itemize}[noitemsep]
  \item Leia código de projetos open-source escritos em C (kernel Linux, Git,
  Redis, SQLite, FFmpeg).
  \item Estude \textit{The C Programming Language} (K\&R) -- o livro original
  de Kernighan e Ritchie.
  \item Aprenda C++ ou Rust -- evoluções diretas de C com segurança e
  abstração modernas.
  \item Estude algoritmos avançados: \textit{Introduction to Algorithms}
  (CLRS), Sedgewick.
  \item Pratique em desafios: Project Euler, Advent of Code, sistemas de
  competição de programação.
\end{itemize}

\textbf{Parabéns por chegar até aqui!} Os 14 capítulos de Deitel \& Deitel
formam um currículo completo e rigoroso. Você está oficialmente pronto para
o próximo nível.

\end{tcolorbox}

\end{document}
